\section{A small invariant for the zero-register jump}
\label{sec:zero-cut}

Let $P$ be the 3,646-site register program of the supplied symmetric cached
machine PA194CS. Its registers are numbered $0,\ldots,8$. Write $X=r_6$.
A decrement instruction $\operatorname{dec}r[z,p]$ goes to $z$ without
changing the register when it is zero, and otherwise subtracts one and
goes to $p$. Define $P_X$ by replacing its instruction at site 26:
\[
 P(26)=\operatorname{jump}34,
 \qquad
 P_X(26)=\operatorname{dec}r_6[34,27].
\]
Every other instruction, including every resolved successor, is unchanged.

The required invariant does not depend on a large abstract-state search.
It uses the following eleven program locations:
\[
 Z_X=\{25,26\}\cup\{3550,3551,3552,3553\}
       \cup\{3633,3634,3635,3636,3637\}.
\]
The entire set of edges entering $Z_X$ from its complement is
\[
\begin{array}{c|c|c}
\text{source}&\text{instruction}&\text{edge into }Z_X\\\hline
22&\operatorname{dec}r_6[25,23]&22\longrightarrow25\\
3545&\operatorname{dec}r_6[3551,3546]&3545\longrightarrow3551\\
3548&\operatorname{dec}r_6[3550,3549]&3548\longrightarrow3550\\
3628&\operatorname{dec}r_6[3634,3629]&3628\longrightarrow3634\\
3631&\operatorname{dec}r_6[3633,3632]&3631\longrightarrow3633
\end{array}
\]
All five are zero branches of a decrement of $r_6$. Inside $Z_X$, the
unchanged instructions are
\[
\begin{array}{r|l@{\qquad}r|l}
25&\operatorname{dec}r_8[27,26]&3550&\operatorname{jump}3553\\
3551&\operatorname{dec}r_1[3553,3552]&3552&\operatorname{jump}3551\\
3553&\operatorname{jump}3637&3633&\operatorname{jump}3637\\
3634&\operatorname{dec}r_4[3637,3635]&3635&\operatorname{inc}r_5\to3636\\
3636&\operatorname{jump}3634&3637&\operatorname{jump}25.
\end{array}
\]
Thus none of these ten instructions changes $r_6$. The eleventh is site
26, whose old instruction preserves it and whose replacement takes its
zero branch whenever $r_6=0$.

\begin{lemma}[Eleven-site zero invariant]
In both $P$ and $P_X$, every run from a configuration satisfying
\[
                    \operatorname{pc}\in Z_X\ \Longrightarrow\ r_6=0
\]
preserves that implication. In particular, it holds from site 25 with
$r_6=0$, irrespective of the values of the other eight registers.
\end{lemma}
\begin{proof}
Induct on the number of register instructions. If the next program counter
is outside $Z_X$, the conclusion is immediate. If the current counter is
in $Z_X$, the displayed instructions preserve a zero value of $r_6$.
If it is outside $Z_X$ and the next counter is inside, one of the five
listed zero branches was taken, so $r_6=0$ at the destination. These cases
exhaust all resolved successors; the accompanying checker verifies that
exhaustiveness against the complete instruction file.
\end{proof}

\begin{theorem}[Exact source-run preservation]
From the clean seed
\[
       (\operatorname{pc},r_0,\ldots,r_8)
                  =(25,2,0,0,0,0,0,0,0,0),
\]
the programs $P$ and $P_X$ have identical register configurations at every
source time and halt at the same source time, if they halt at all. More
generally this holds from every configuration satisfying the preceding
invariant.
\end{theorem}
\begin{proof}
Induct simultaneously on the two runs. The instructions agree except at
site 26. At that site the invariant gives $r_6=0$, and both instructions
preserve every register and go to site 34. Thus the complete next
configurations agree, not merely their halting outcomes. Their identical
halt instruction is preserved. The invariant and the run induction are
ordinary arithmetical inductions on coded finite runs, so PA proves the
corresponding universal statements. No consistency or soundness assumption
about PA is used.
\end{proof}

The binary-table modification implementing this source replacement is
native row $(27,1):0L40\mapsto1R21$. The source proof alone does not justify
that row edit: it is separately bound to the exact table by complete
symbolic register contracts, including the now-unreachable positive
branch at site 26, and by complete-tape blank initialization. The two
binary machines need not have equal physical running times or equal
spare lists. Their own positive-duration simulations to the respective
source programs, followed by the exact source theorem, are sufficient.

\paragraph{An alternative using the command register.}
Replacing site 26 instead by $\operatorname{dec}r_3[34,27]$ has the same
source-run preservation property. One suitable zero cut is
\[
       Z_A=\{0,\ldots,9\}\cup\{11,\ldots,26\}
                 \cup\{3523,\ldots,3645\}.
\]
It contains 149 sites. Every edge into it from outside is a zero branch
of a decrement of $r_3$, and its internal instructions preserve $r_3=0$.
The same induction applies. The smaller $Z_X$ was chosen for the primary
proof because its entire boundary can be printed in five rows.
